% GENERATED FILE. Edit canonical lessons, then run npm run generate:books.
% canonical-source-hash: fnv1a64:9a28e257cc0d444a
% canonical-lessons: KA-C287-dada, KA-C287-bara, KA-C287-bhukampa, KA-C287-prapanca, KA-C287-gode

\chapter{Riverbank, Drought, Earthquake, World, Wall}
\label{ch:ka-riverbank-drought-earthquake-world-wall}
\hlchaptermodality{287}

\begin{chapteropening}
\textbf{By the end of this chapter:} \emph{I can say a riverbank, a drought, an earthquake, the world and a wall.}

Complete the last lesson of chapter 287: I can say a riverbank, a drought, an earthquake, the world and a wall.
\end{chapteropening}

\section[\texorpdfstring{\kn{ದಡ} (daḍa)}{ದಡ (daḍa)}]{\kn{ದಡ} (daḍa) — a riverbank}

\label{lesson:KA-C287-dada}



\begin{quote}
Before the new one: say the Kannada for a platform, then the Kannada for abroad.
\end{quote}

\subsection*{You'll want to know: \kn{ದಡ}}
\textbf{\kn{ದಡ}} — \emph{daḍa} — \textquotedblleft{}a riverbank\textquotedblright{}.

\begin{cousinweb}
\textbf{\kn{ದಡ}} (\emph{daḍa}) is probably a tadbhava of Sanskrit \emph{taṭa}, \textquotedblleft{}slope, shore\textquotedblright{} (CDIAL 5629), through Prakrit \emph{taḍa}; CDIAL does not list the Kannada word.
\end{cousinweb}

\begin{grammarlens}[title={What you've built}]
One more word for this chapter.
\end{grammarlens}

\subsection*{Your turn}
\begin{itemize}
\raggedright
  \item \emph{Say it:} \emph{daḍa}
  \item \emph{Say it:} \emph{daḍa}, once more
  \item \emph{Say it:} say \emph{daḍa}
  \item \emph{Recall:} say the Kannada for a platform, then the Kannada for abroad, then say \emph{daḍa} again
\end{itemize}

\begin{tcolorbox}[breakable,colback=teal!4,colframe=teal!35!black,title={Before you move on}]
What does \kn{ದಡ} mean? (\textquotedblleft{}A riverbank\textquotedblright{}.) Say it once more.
\end{tcolorbox}
\section[\texorpdfstring{\kn{ಬರ} (bara)}{ಬರ (bara)}]{\kn{ಬರ} (bara) — a drought}

\label{lesson:KA-C287-bara}



\begin{quote}
Before the new one: say the Kannada for abroad, then the Kannada for a riverbank.
\end{quote}

\subsection*{You'll want to know: \kn{ಬರ}}
\textbf{\kn{ಬರ}} — \emph{bara} — \textquotedblleft{}a drought\textquotedblright{}.

\begin{cousinweb}
\textbf{\kn{ಬರ}} (\emph{bara}) is native Dravidian, one word with Tamil \emph{vaṟa}, Telugu \emph{vaṟapu} and Malayalam \emph{vaṟu} (DEDR 5320). DEDR lists older \emph{baṟa}, \textquotedblleft{}dryness, barrenness, drought\textquotedblright{}.
\end{cousinweb}

\begin{grammarlens}[title={What you've built}]
One more word for this chapter.
\end{grammarlens}

\subsection*{Your turn}
\begin{itemize}
\raggedright
  \item \emph{Say it:} \emph{bara}
  \item \emph{Say it:} \emph{bara}, once more
  \item \emph{Say it:} say \emph{bara}
  \item \emph{Recall:} say the Kannada for abroad, then the Kannada for a riverbank, then say \emph{bara} again
\end{itemize}

\begin{tcolorbox}[breakable,colback=teal!4,colframe=teal!35!black,title={Before you move on}]
What does \kn{ಬರ} mean? (\textquotedblleft{}A drought\textquotedblright{}.) Say it once more.
\end{tcolorbox}
\section[\texorpdfstring{\kn{ಭೂಕಂಪ} (bhūkampa)}{ಭೂಕಂಪ (bhūkampa)}]{\kn{ಭೂಕಂಪ} (bhūkampa) — an earthquake}

\label{lesson:KA-C287-bhukampa}



\begin{quote}
Before the new one: say the Kannada for a riverbank, then the Kannada for a drought.
\end{quote}

\subsection*{You'll want to know: \kn{ಭೂಕಂಪ}}
\textbf{\kn{ಭೂಕಂಪ}} — \emph{bhūkampa} — \textquotedblleft{}an earthquake\textquotedblright{}.

\begin{cousinweb}
\textbf{\kn{ಭೂಕಂಪ}} (\emph{bhūkampa}) is Sanskrit \emph{bhūkampa}, \textquotedblleft{}earthquake\textquotedblright{}, taken over whole as a learned loan, a tatsama (CDIAL 9547). \emph{Bhū}, \textquotedblleft{}earth\textquotedblright{}, and \emph{kampa}, \textquotedblleft{}trembling\textquotedblright{}.
\end{cousinweb}

\begin{grammarlens}[title={What you've built}]
One more word for this chapter.
\end{grammarlens}

\subsection*{Your turn}
\begin{itemize}
\raggedright
  \item \emph{Say it:} \emph{bhūkampa}
  \item \emph{Say it:} \emph{bhūkampa}, once more
  \item \emph{Say it:} say \emph{bhūkampa}
  \item \emph{Recall:} say the Kannada for a riverbank, then the Kannada for a drought, then say \emph{bhūkampa} again
\end{itemize}

\begin{tcolorbox}[breakable,colback=teal!4,colframe=teal!35!black,title={Before you move on}]
What does \kn{ಭೂಕಂಪ} mean? (\textquotedblleft{}An earthquake\textquotedblright{}.) Say it once more.
\end{tcolorbox}
\section[\texorpdfstring{\kn{ಪ್ರಪಂಚ} (prapañca)}{ಪ್ರಪಂಚ (prapañca)}]{\kn{ಪ್ರಪಂಚ} (prapañca) — the world}

\label{lesson:KA-C287-prapanca}



\begin{quote}
Before the new one: say the Kannada for a drought, then the Kannada for an earthquake.
\end{quote}

\subsection*{You'll want to know: \kn{ಪ್ರಪಂಚ}}
\textbf{\kn{ಪ್ರಪಂಚ}} — \emph{prapañca} — \textquotedblleft{}the world\textquotedblright{}.

\begin{cousinweb}
\textbf{\kn{ಪ್ರಪಂಚ}} (\emph{prapañca}) is Sanskrit \emph{prapañca}, \textquotedblleft{}expansion, the visible world\textquotedblright{}.
\end{cousinweb}

\begin{grammarlens}[title={What you've built}]
One more word for this chapter.
\end{grammarlens}

\subsection*{Your turn}
\begin{itemize}
\raggedright
  \item \emph{Say it:} \emph{prapañca}
  \item \emph{Say it:} \emph{prapañca}, once more
  \item \emph{Say it:} say \emph{prapañca}
  \item \emph{Recall:} say the Kannada for a drought, then the Kannada for an earthquake, then say \emph{prapañca} again
\end{itemize}

\begin{tcolorbox}[breakable,colback=teal!4,colframe=teal!35!black,title={Before you move on}]
What does \kn{ಪ್ರಪಂಚ} mean? (\textquotedblleft{}The world\textquotedblright{}.) Say it once more.
\end{tcolorbox}
\section[\texorpdfstring{\kn{ಗೋಡೆ} (gōḍe)}{ಗೋಡೆ (gōḍe)}]{\kn{ಗೋಡೆ} (gōḍe) — a wall}

\label{lesson:KA-C287-gode}



\begin{quote}
Before the new one: say the Kannada for an earthquake, then the Kannada for the world.
\end{quote}

\subsection*{You'll want to know: \kn{ಗೋಡೆ}}
\textbf{\kn{ಗೋಡೆ}} — \emph{gōḍe} — \textquotedblleft{}a wall\textquotedblright{}.

\begin{cousinweb}
\textbf{\kn{ಗೋಡೆ}} (\emph{gōḍe}) is native Dravidian, one word with Telugu \emph{kōṭa} (DEDR 2207). DEDR lists \emph{gōḍe} beside \emph{kōṭe}, \textquotedblleft{}fort, rampart\textquotedblright{}. DEDR 2207 compares Sanskrit \emph{koṭṭa}.
\end{cousinweb}

\begin{grammarlens}[title={What you've built}]
One more word for this chapter.
\end{grammarlens}

\subsection*{Your turn}
\begin{itemize}
\raggedright
  \item \emph{Say it:} \emph{gōḍe}
  \item \emph{Say it:} \emph{gōḍe}, once more
  \item \emph{Say it:} say \emph{gōḍe}
  \item \emph{Recall:} say the Kannada for an earthquake, then the Kannada for the world, then say \emph{gōḍe} again
\end{itemize}

\begin{tcolorbox}[breakable,colback=teal!4,colframe=teal!35!black,title={Before you move on}]
What does \kn{ಗೋಡೆ} mean? (\textquotedblleft{}A wall\textquotedblright{}.) Say it once more.
\end{tcolorbox}
